\section{Introduction}
A cluster operator sees tasks belonging to jobs and machines; an incident
manager sees unresolved tickets linked to handlers and teams. Both need
collective forecasts: how many current entities
will remain unresolved, and which groups might fall behind together? Temporal
models can supply each entity's event probability, but those probabilities
alone do not specify the distribution of group counts or the largest group
delay. Moreover, a well-scored count distribution need not support the
particular decision quantile used by an operator.

This study asks a practical evaluation question: \emph{what additional evidence
is needed before individual forecasts on a temporal operational graph can be
used as collective scenarios?} Three issues are easily conflated. Individual
predictions may share a mean bias. Outcomes may retain common variation after
conditioning on observed history. Typed relations may identify additional
association beyond a global factor. Improving the first issue, adding the
second, and exploiting the third require different controls. An apparent
``graph benefit'' against independent sampling can otherwise be mostly a
benefit of shared uncertainty available without typed relations.

We define a common current-cohort target for six public logs: a compute-cluster
trace, Azure VMs, two loan-workflow releases, and two IT-incident sources.
At each origin, a timestamp-bounded
adapter reconstructs the active entities and their observed memberships.
Scenarios assign each entity a first-exit bin across four horizons, including
survival beyond the final horizon. The cohort is not selected using future
exits, arrivals or assignments. The resulting counts can be projected onto
overlapping jobs and machines or onto recent-handler groups.

The contribution is a reproducible benchmark and controlled empirical study,
not a new copula architecture. We compare strong individual predictors,
matched-marginal global and typed factors, direct count distributions,
historical residuals, and discrete probabilistic reconciliation. Delayed
calibration controls separate mean error from dependence, and threshold-free
joint scores test whether count accuracy extends to a group-risk profile.
An initial two-log evaluation motivates retrospective diagnostics; a separately
frozen third-log period and three further frozen source periods supply
confirmation. Storage-bounded transportation pilots and incomplete-history
audits are retained as negative or deferred evidence. The manuscript identifies
these stages explicitly rather than pooling development and confirmation.

Three findings organize the study. Shared factors are strong controls, so the
incremental contribution of typed incidence must be measured against them.
Job grouping and IT handler/team incidence add predictive signal in particular
settings, but graph architecture, relation identity, and dependence are not
interchangeable sources of improvement. Calibration can materially change
both collective scores and asymmetric loss,
but transferring residuals across predictor versions is unreliable. Finally,
methods that forecast individual group counts well can rank differently on
the largest standardized group delay. These findings motivate an evaluation
workflow that checks marginal bias, group-count uncertainty, and joint group
risk separately, rather than treating one score as validation of a scenario
generator.
